\documentclass[10pt,a4paper]{amsart}
\usepackage{amsmath,amssymb,amsthm,graphicx}
\usepackage{verbatim}
\newenvironment{reference}{\comment}{\endcomment}
%\newenvironment{reference}{}{}
\newenvironment{slogan}{\comment}{\endcomment}
\newenvironment{history}{\comment}{\endcomment}

% For commutative diagrams we use Xy-pic
\usepackage[all]{xy}

% We use 2cell for 2-commutative diagrams.
\xyoption{2cell}
% Theorem environments.
%
\theoremstyle{plain}
\newtheorem{theorem}[subsection]{Theorem}
\newtheorem{proposition}[subsection]{Proposition}
\newtheorem{lemma}[subsection]{Lemma}

\theoremstyle{definition}
\newtheorem{definition}[subsection]{Definition}
\newtheorem{example}[subsection]{Example}
\newtheorem{exercise}[subsection]{Exercise}
\newtheorem{situation}[subsection]{Situation}

\theoremstyle{remark}
\newtheorem{remark}[subsection]{Remark}
\newtheorem{remarks}[subsection]{Remarks}

\numberwithin{equation}{subsection}

% Macros
%
\def\lim{\mathop{\mathrm{lim}}\nolimits}
\def\colim{\mathop{\mathrm{colim}}\nolimits}
\def\Spec{\mathop{\mathrm{Spec}}}
\def\Hom{\mathop{\mathrm{Hom}}\nolimits}
\def\Ext{\mathop{\mathrm{Ext}}\nolimits}
\def\SheafHom{\mathop{\mathcal{H}\!\mathit{om}}\nolimits}
\def\SheafExt{\mathop{\mathcal{E}\!\mathit{xt}}\nolimits}
\def\Sch{\mathit{Sch}}
\def\Mor{\mathop{\mathrm{Mor}}\nolimits}
\def\Ob{\mathop{\mathrm{Ob}}\nolimits}
\def\Sh{\mathop{\mathit{Sh}}\nolimits}
\def\NL{\mathop{N\!L}\nolimits}
\def\CH{\mathop{\mathrm{CH}}\nolimits}
\def\proetale{{pro\text{-}\acute{e}tale}}
\def\etale{{\acute{e}tale}}
\def\QCoh{\mathit{QCoh}}
\def\Ker{\mathop{\mathrm{Ker}}}
\def\Im{\mathop{\mathrm{Im}}}
\def\Coker{\mathop{\mathrm{Coker}}}
\def\Coim{\mathop{\mathrm{Coim}}}

% Boxtimes
%
\DeclareMathSymbol{\boxtimes}{\mathbin}{AMSa}{"02}

%
% Macros for moduli stacks/spaces
%
\def\QCohstack{\mathcal{QC}\!\mathit{oh}}
\def\Cohstack{\mathcal{C}\!\mathit{oh}}
\def\Spacesstack{\mathcal{S}\!\mathit{paces}}
\def\Quotfunctor{\mathrm{Quot}}
\def\Hilbfunctor{\mathrm{Hilb}}
\def\Curvesstack{\mathcal{C}\!\mathit{urves}}
\def\Polarizedstack{\mathcal{P}\!\mathit{olarized}}
\def\Complexesstack{\mathcal{C}\!\mathit{omplexes}}
% \Pic is the operator that assigns to X its picard group, usage \Pic(X)
% \Picardstack_{X/B} denotes the Picard stack of X over B
% \Picardfunctor_{X/B} denotes the Picard functor of X over B
\def\Pic{\mathop{\mathrm{Pic}}\nolimits}
\def\Picardstack{\mathcal{P}\!\mathit{ic}}
\def\Picardfunctor{\mathrm{Pic}}
\def\Deformationcategory{\mathcal{D}\!\mathit{ef}}
\usepackage{cleveref}
\crefname{originallemma}{original lemma}{original lemmas}
\Crefname{originallemma}{Original lemma}{Original lemmas}
\makeatletter
\newlabel{algebra-lemma-apply-grothendieck}{{0337}{source}{}{sourcepointer-1}{}}
\newlabel{algebra-lemma-apply-grothendieck@cref}{{[originallemma][0][]0337}{[1][0][]source}{}{}{}}
\makeatother
\crefname{originallemma}{original lemma}{original lemmas}
\Crefname{originallemma}{Original lemma}{Original lemmas}
\makeatletter
\newlabel{algebra-lemma-completion-flat}{{00MB}{source}{}{sourcepointer-2}{}}
\newlabel{algebra-lemma-completion-flat@cref}{{[originallemma][0][]00MB}{[1][0][]source}{}{}{}}
\makeatother
\crefname{originallemma}{original lemma}{original lemmas}
\Crefname{originallemma}{Original lemma}{Original lemmas}
\makeatletter
\newlabel{algebra-lemma-criterion-normal}{{031S}{source}{}{sourcepointer-3}{}}
\newlabel{algebra-lemma-criterion-normal@cref}{{[originallemma][0][]031S}{[1][0][]source}{}{}{}}
\makeatother
\crefname{originallemma}{original lemma}{original lemmas}
\Crefname{originallemma}{Original lemma}{Original lemmas}
\makeatletter
\newlabel{algebra-lemma-descent-Rk}{{0353}{source}{}{sourcepointer-4}{}}
\newlabel{algebra-lemma-descent-Rk@cref}{{[originallemma][0][]0353}{[1][0][]source}{}{}{}}
\makeatother
\crefname{originallemma}{original lemma}{original lemmas}
\Crefname{originallemma}{Original lemma}{Original lemmas}
\makeatletter
\newlabel{algebra-lemma-descent-Sk}{{0352}{source}{}{sourcepointer-5}{}}
\newlabel{algebra-lemma-descent-Sk@cref}{{[originallemma][0][]0352}{[1][0][]source}{}{}{}}
\makeatother
\crefname{originallemma}{original lemma}{original lemmas}
\Crefname{originallemma}{Original lemma}{Original lemmas}
\makeatletter
\newlabel{algebra-lemma-dimension-base-fibre-total}{{00OM}{source}{}{sourcepointer-6}{}}
\newlabel{algebra-lemma-dimension-base-fibre-total@cref}{{[originallemma][0][]00OM}{[1][0][]source}{}{}{}}
\makeatother
\title{Normality of a subalgebra from normality of the algebra and of its completion}
\author{The Stacks Project Authors}
\date{stacks-project:resolve.tex:local-snapshot-20261012; GNU FDL 1.2 or later}
\begin{document}
\maketitle
\noindent\textit{Editorial note.} All delivered mathematical text is the author's own, copied from the pinned Stacks Project chapter file resolve.tex; no statement or proof step was rewritten, reordered or removed.\par\medskip
\noindent\textit{Editorial note.} The counted proof is the author's complete argument: the reduction through Serre's criterion and the exactness of completion, the verification of (S2) split into the higher and lower dimensional cases using depth and faithful flatness, and the verification of (R1) through the dimension formula and descent along a faithfully flat map.\par\medskip
\noindent\textit{Editorial note.} The label lemma-algebra-helper names the counted statement, and the proof environment immediately following it in the source is delivered verbatim as its proof.\par\medskip
\noindent\textit{Editorial note.} The counted statement is carried by the author's own lemma environment in the chapter source.\par\medskip
\noindent\textit{Editorial note.} Pre-open gate: python operations/agent-authoring/tools/precard\_gate.py --source resolve.tex --label lemma-algebra-helper --start 3775 --end 3824 --local-id 28092 returned RESULT=OK (L28092 free, occupancy FREE, public interval clear). Pre-assembly tag self-check: canonical tag resolved and not present in reference/public-items.jsonl.\par\medskip
\noindent\textit{Editorial note.} The statement and proof are copied from the pinned chapter file; the cited inputs are kept as external pointers with their canonical Stacks tags.\par\medskip
\noindent\textit{Allowed standard prerequisites (editorial).}\begin{itemize}
\item Serre's criterion for normality. Algebra, the lemma characterizing a normal Noetherian ring by the conditions (R1) and (S2) on its local rings. Supplies the two conditions verified in the proof.
\item Flatness and exactness of the adic completion, and depth behaviour under flat maps. Algebra, the lemma that completion is given by tensoring with a flat algebra, and the lemma controlling the depth of a flat local extension. Supplies the exactness of completion on finite modules and the depth estimate used in the high dimensional case.
\item Descent of the conditions (Rk) and (Sk) along faithfully flat maps. Algebra, the lemmas descending the Serre conditions along a faithfully flat ring map. Supply the conclusions that the local rings of the completion satisfy (S2), respectively (R1).
\end{itemize}\medskip
\ifdefined\hypertarget\hypertarget{sourcepointer-1}{}\fi
\noindent\textit{Source reference.} Original lemma 0337: sources/stacks-chow/tags.txt, source lines 4046--4046. Prerequisite pointer invoked at proof line(s) 3803; it names a standard fact of a sibling chapter that the counted argument uses, and every step of the counted proof built on it is retained verbatim in the delivered body.\par\medskip
\ifdefined\hypertarget\hypertarget{sourcepointer-2}{}\fi
\noindent\textit{Source reference.} Original lemma 00MB: sources/stacks-chow/tags.txt, source lines 817--817. Prerequisite pointer invoked at proof line(s) 3794; it names a standard fact of a sibling chapter that the counted argument uses, and every step of the counted proof built on it is retained verbatim in the delivered body.\par\medskip
\ifdefined\hypertarget\hypertarget{sourcepointer-3}{}\fi
\noindent\textit{Source reference.} Original lemma 031S: sources/stacks-chow/tags.txt, source lines 3995--3995. Prerequisite pointer invoked at proof line(s) 3796; it names a standard fact of a sibling chapter that the counted argument uses, and every step of the counted proof built on it is retained verbatim in the delivered body.\par\medskip
\ifdefined\hypertarget\hypertarget{sourcepointer-4}{}\fi
\noindent\textit{Source reference.} Original lemma 0353: sources/stacks-chow/tags.txt, source lines 4114--4114. Prerequisite pointer invoked at proof line(s) 3823; it names a standard fact of a sibling chapter that the counted argument uses, and every step of the counted proof built on it is retained verbatim in the delivered body.\par\medskip
\ifdefined\hypertarget\hypertarget{sourcepointer-5}{}\fi
\noindent\textit{Source reference.} Original lemma 0352: sources/stacks-chow/tags.txt, source lines 4113--4113. Prerequisite pointer invoked at proof line(s) 3812; it names a standard fact of a sibling chapter that the counted argument uses, and every step of the counted proof built on it is retained verbatim in the delivered body.\par\medskip
\ifdefined\hypertarget\hypertarget{sourcepointer-6}{}\fi
\noindent\textit{Source reference.} Original lemma 00OM: sources/stacks-chow/tags.txt, source lines 900--900. Prerequisite pointer invoked at proof line(s) 3819; it names a standard fact of a sibling chapter that the counted argument uses, and every step of the counted proof built on it is retained verbatim in the delivered body.\par\medskip
\begin{lemma}
\label{lemma-algebra-helper}
Let $(A, \mathfrak m)$ be a local Noetherian ring. Let $B \subset C$
be finite $A$-algebras. Assume that (a) $B$ is a normal ring, and
(b) the $\mathfrak m$-adic completion $C^\wedge$ is a normal ring.
Then $B^\wedge$ is a normal ring.
\end{lemma}

\begin{proof}
Consider the commutative diagram
$$
\xymatrix{
B \ar[r] \ar[d] & C \ar[d] \\
B^\wedge \ar[r] & C^\wedge
}
$$
Recall that $\mathfrak m$-adic completion on the category of
finite $A$-modules is exact because it is given by tensoring with
the flat $A$-algebra $A^\wedge$
(Algebra, Lemma \ref{algebra-lemma-completion-flat}).
We will use Serre's criterion
(Algebra, Lemma \ref{algebra-lemma-criterion-normal})
to prove that the Noetherian ring $B^\wedge$ is normal.
Let $\mathfrak q \subset B^\wedge$ be a prime lying over
$\mathfrak p \subset B$. If $\dim(B_\mathfrak p) \geq 2$, then
$\text{depth}(B_\mathfrak p) \geq 2$ and since
$B_\mathfrak p \to B^\wedge_\mathfrak q$ is flat we find
that $\text{depth}(B^\wedge_\mathfrak q) \geq 2$
(Algebra, Lemma \ref{algebra-lemma-apply-grothendieck}).
If $\dim(B_\mathfrak p) \leq 1$, then $B_\mathfrak p$ is
either a discrete valuation ring or a field.
In that case $C_\mathfrak p$ is faithfully flat over $B_\mathfrak p$
(because it is finite and torsion free).
Hence $B^\wedge_\mathfrak p \to C^\wedge_\mathfrak p$ is
faithfully flat and the same holds after localizing at $\mathfrak q$.
As $C^\wedge$ and hence any localization is $(S_2)$ we conclude that
$B^\wedge_\mathfrak p$ is $(S_2)$ by
Algebra, Lemma \ref{algebra-lemma-descent-Sk}.
All in all we find that
$(S_2)$ holds for $B^\wedge$. To prove that $B^\wedge$ is
$(R_1)$ we only have to consider primes $\mathfrak q \subset B^\wedge$
with $\dim(B^\wedge_\mathfrak q) \leq 1$. Since
$\dim(B^\wedge_\mathfrak q) = \dim(B_\mathfrak p) +
\dim(B^\wedge_\mathfrak q/\mathfrak p B^\wedge_\mathfrak q)$ by
Algebra, Lemma \ref{algebra-lemma-dimension-base-fibre-total}
we find that $\dim(B_\mathfrak p) \leq 1$ and
we see that $B^\wedge_\mathfrak q \to C^\wedge_\mathfrak q$
is faithfully flat as before. We conclude using
Algebra, Lemma \ref{algebra-lemma-descent-Rk}.
\end{proof}

\end{document}
